%% Appendix, for the supplementary material (outside the 20-page limit).
%% Holds the formal definitions that the paper uses but does not discuss
%% in detail (rule of thumb: what the text discusses in detail is defined
%% in the paper; the rest goes here).

\appendix

\section{Type Systems}
\label{app:types}

%% From definitions.tex (Types and Languages) and agda/Mu/Language/*.agda.
Every language of Section~\ref{sec:compiler} is intrinsically typed in
the Agda: a term is indexed by its context and its type, so only well
typed terms can be written. This appendix gives the typing rules in the
notation of Figure~\ref{fig:grammars}. The base types are
$\kw{Nat}$, $\kw{Int}$, $\kw{Bool}$, $\kw{Unit}$, and $\kw{Void}$, with
$\kw{Bool} = \kw{tup}() + \kw{tup}()$. The result type of a primitive
is $\kw{Int}$ for $\kw{sub}$ and $\kw{Bool}$ for $\kw{eq}$. The body of
a recursive type $\mu X.\, A$ is contractive: it starts with a base
type, $\rightarrow$, $\kw{tup}$, $+$, or $\kw{code}$.

\subsection{\iswim}
A judgment $\Gamma \vdash M : A$ types a term in a context $\Gamma$ of
variable types.
\[
\begin{array}{@{}r@{\;:\;}l@{\qquad}l@{}}
  x & A & x{:}A \in \Gamma \\
  \lambda x.\, N & A \rightarrow C & \Gamma, x{:}A \vdash N : C \\
  L\;M & C & L : A \rightarrow C,\ M : A \\
  c & B & c \text{ a constant of base type } B \\
  \kw{tuple}\;(M_1, \ldots, M_n) & \kw{tup}(A_1, \ldots, A_n) & M_i : A_i \\
  \kw{get}_i\,M & A_i & M : \kw{tup}(A_1, \ldots, A_n),\ 1 \leq i \leq n \\
  \kw{inl}\;M,\ \kw{inr}\;M & A + C & M : A, \text{ resp. } M : C \\
  \kw{case}\;L\;(x.\,M)\;(y.\,N) & D & L : A + C,\ \Gamma, x{:}A \vdash M : D,\ \Gamma, y{:}C \vdash N : D \\
  \kw{prim}_p\;L\;M & \mathit{out}(p) & L, M : \kw{Int} \\
  \kw{let}\;x = M\;\kw{in}\;N & C & M : A,\ \Gamma, x{:}A \vdash N : C \\
  \kw{fold}\;M & \mu X.\, A & M : A[\mu X.\, A / X] \\
  \kw{unfold}\;M & A[\mu X.\, A / X] & M : \mu X.\, A
\end{array}
\]
Each line reads: the term on the left has the type in the middle, in
$\Gamma$, provided the premises on the right hold, with the subterms
typed in $\Gamma$ unless stated otherwise.

\subsection{The Closure Languages}
Clos, EnvClos, FlatClos, and GlobalFun type their other forms as \iswim\
does. A closure's type does not show the types $\Delta = A_1, \ldots,
A_n$ of its captured values.
\[
  \frac{\Gamma \vdash M_i : A_i \ (1 \leq i \leq n) \qquad
        x_1{:}A_1, \ldots, x_n{:}A_n, y{:}A \vdash N : C}
       {\Gamma \vdash \kw{clos}\;N\;\vec{M} : A \rightarrow C}\;\text{(Clos)}
\]
\[
  \frac{\Gamma \vdash M_i : A_i \ (1 \leq i \leq n) \qquad
        e{:}\kw{tup}(\Delta),\, y{:}A \vdash N : C}
       {\Gamma \vdash \kw{clos}\;N\;\vec{M} : A \rightarrow C}\;\text{(EnvClos)}
\]
FlatClos types $\kw{code}$, $\kw{pack}$, and $\kw{callc}$ as in
Section~\ref{sec:compiler}. In GlobalFun, a program has a table
$\Sigma = s_0, \ldots, s_{m-1}$ of the signatures
$s_k = (\Delta_k; A_k \rightarrow C_k)$ of its definitions, and a term
is typed by the table as well as its context, $\Sigma; \Gamma \vdash M
: A$:
\[
  \frac{k < m}{\Sigma; \Gamma \vdash f_k : \kw{code}(\Delta_k; A_k \rightarrow C_k)}
\]
\[
  \frac{\Sigma; e{:}\kw{tup}(\Delta_k), y{:}A_k \vdash N_k : C_k \ (d_k = \kw{code}\;N_k,\ k < m)
        \quad \Sigma; \emptyset \vdash M : A}
       {\vdash (d_0, \ldots, d_{m-1};\; M) : A}
\]
Mon is typed as GlobalFun, with atoms in place of operands.

\subsection{Imp}
A body runs in a frame of slots, typed by the list
$\Delta = A_0, \ldots, A_{j-1}$ of the types of the slots assigned so
far; slot $s < j$ has type $A_s$. An assignment and a $\kw{case}$ write
the next slot, $|\Delta|$.
\[
  \Sigma; \Delta \vdash s : A_s
  \qquad
  \frac{\Sigma; \Delta \vdash e_1 : A \quad \Sigma; \Delta, A \vdash e_2 : C}
       {\Sigma; \Delta \vdash |\Delta| \mathrel{:=} e_1;\; e_2 : C}
\]
\[
  \frac{\Sigma; \Delta \vdash a : A + B \quad
        \Sigma; \Delta, A \vdash e_1 : C \quad \Sigma; \Delta, B \vdash e_2 : C}
       {\Sigma; \Delta \vdash \kw{case}\;a\;(|\Delta|.\,e_1)\;(|\Delta|.\,e_2) : C}
\]
A definition $e_k$ is typed in the frame $\kw{tup}(\Delta_k), A_k$ (its
environment in slot 0 and its argument in slot 1) at type $C_k$, and
the main expression in the empty frame.

\subsection{CFG}
A tail is typed by the type $R$ of the body's result, the list
$\vec{\Delta} = \Delta_0, \ldots, \Delta_{j-1}$ of the frames of the
blocks it may jump to, and its own frame $\Delta$. A label $l$ is an
index into $\vec{\Delta}$. Block $t_l$ is typed over the blocks older
than it, so control flow is acyclic. Here $\Delta' \sqsubseteq \Delta$
says that $\Delta'$ is a prefix of $\Delta$; an assignment to slot
$|\Delta'|$ forgets the slots from $|\Delta'|$ on.
\[
  \frac{\Sigma; \Delta \vdash r : R}{\Sigma; R; \vec{\Delta}; \Delta \vdash \kw{return}\;r}
  \qquad
  \frac{\Delta' \sqsubseteq \Delta \quad \Sigma; \Delta \vdash r : A \quad
        \Sigma; R; \vec{\Delta}; \Delta', A \vdash t}
       {\Sigma; R; \vec{\Delta}; \Delta \vdash |\Delta'| \mathrel{:=} r;\; t}
  \qquad
  \frac{\Delta_l = \Delta}{\Sigma; R; \vec{\Delta}; \Delta \vdash \kw{goto}\;l}
\]
\[
  \frac{\Sigma; \Delta \vdash a : A + B \quad \Delta' \sqsubseteq \Delta \quad
        \Delta_{l_1} = \Delta', A \quad \Delta_{l_2} = \Delta', B}
       {\Sigma; R; \vec{\Delta}; \Delta \vdash
        \kw{case}\;a\;(|\Delta'|.\,l_1)\;(|\Delta'|.\,l_2)}
\]
\[
  \frac{\Sigma; R; \Delta_0, \ldots, \Delta_{l-1}; \Delta_l \vdash t_l \ (l < j) \quad
        \Sigma; R; \vec{\Delta}; \Delta \vdash t}
       {\Sigma; R; \Delta \vdash (t_0, \ldots, t_{j-1};\; t)}
\]

\subsection{X86$\mu$}
An instruction is typed by the type of the machine state it starts in
and the one it leaves. A state type
$\psi = \langle \Theta; \Delta; O; F \rangle$ gives the types $\Theta$
of the registers, the frame $\Delta$ of the slots assigned so far, the
open closure $O$ (none, written $\cdot$, or a slot $x$ of $\Delta$ of a
function type), and whether the flag $F$ is set. A register holds a
value of a type $A$, the code $\kw{code}_\alpha(A \rightarrow C)$ or the
environment $\kw{tup}(\alpha)$ of the open closure, or nothing. The
type variable $\alpha$ stands for the hidden types of the environment
of the open closure. An operand ${\uparrow}a$ has type $\mu X.\, A$ when
$a$ has type $A[\mu X.\, A / X]$, and ${\downarrow}a$ the reverse. A
destination $d$ is written at a type $\tau$,
$n \vdash d : \psi \Rightarrow_\tau \psi'$: a register write sets the
register's type, and a write to slot $|\Delta'|$, for a prefix
$\Delta'$ of $\Delta$ with $|\Delta'| < n$, leaves the frame
$\Delta', A$ and closes the open closure. The instructions are typed
as follows, with $\psi = \langle \Theta; \Delta; O; F \rangle$:
\[
\footnotesize
\begin{array}{@{}l@{\;:\;}l@{\quad}l@{}}
  \kw{movq}\; a, d & \psi \Rightarrow \psi' & a : \tau,\ d : \psi \Rightarrow_\tau \psi' \\
  \kw{subq}\; a, \%r & \psi \Rightarrow \psi & a : \kw{Int},\ \Theta(r) = \kw{Int} \\
  \kw{cmpq}\; a_2, a_1 & \psi \Rightarrow \langle \Theta; \Delta; O; \kw{set} \rangle & a_1, a_2 : \kw{Int} \\
  \kw{sete}\; d & \langle \Theta; \Delta; O; \kw{set} \rangle \Rightarrow \psi' & d : \langle \Theta; \Delta; O; \kw{set} \rangle \Rightarrow_{\kw{Bool}} \psi' \\
  \kw{alloc}\; o, d & \psi \Rightarrow \psi' & o : A,\ d : \psi \Rightarrow_A \psi' \\
  \kw{load}\; \kw{get}_i\; a, d & \psi \Rightarrow \psi' & a : \kw{tup}(A_1, \ldots, A_m),\ d : \psi \Rightarrow_{A_i} \psi' \\
  \kw{load}\; \kw{fst}\; x, \%r & \psi \Rightarrow \langle \mathit{wipe}(\Theta)[r{:}\kw{code}_\alpha(A \rightarrow C)]; \Delta; x; F \rangle & x : A \rightarrow C \in \Delta \\
  \kw{load}\; \kw{snd}\; x, \%r & \langle \Theta; \Delta; x; F \rangle \Rightarrow \langle \Theta[r{:}\kw{tup}(\alpha)]; \Delta; x; F \rangle & \\
  \kw{callq}\; a & \psi \Rightarrow \langle \kw{rax}{:}C; \Delta; \cdot; \kw{unset} \rangle
    & a : \kw{code}_\alpha(A \rightarrow C),\ \Theta(\kw{rdi}) = \kw{tup}(\alpha), \\
  \multicolumn{2}{l}{} & \Theta(\kw{rsi}) = A
\end{array}
\]
Here $\mathit{wipe}(\Theta)$ unsets the registers that hold the code or
the environment of an earlier open closure, so $\alpha$ always stands
for the environment of one closure. An object $\kw{pair}(a_1, a_2)$ has
type $A \rightarrow C$ when $a_1 : \kw{code}(\Delta; A \rightarrow C)$
and $a_2 : \kw{tup}(\Delta)$; tuples and injections are typed as in
\iswim. A sequence $\beta$ is typed by the result type $R$, the frames
$\vec{\Delta}$ of the blocks it may jump to, and its starting state:
\[
  \frac{\Sigma; n \vdash i : \psi \Rightarrow \psi' \quad \Sigma; n; R; \vec{\Delta} \vdash \beta : \psi'}
       {\Sigma; n; R; \vec{\Delta} \vdash i;\; \beta : \psi}
\]
\[
  \frac{\Theta(\kw{rax}) = R}{\Sigma; n; R; \vec{\Delta} \vdash \kw{retq} : \langle \Theta; \Delta; O; F \rangle}
  \qquad
  \frac{\Delta_l = \Delta}{\Sigma; n; R; \vec{\Delta} \vdash \kw{jmp}\;l : \langle \Theta; \Delta; O; F \rangle}
\]
\[
  \frac{x : A + B \in \Delta \quad \Delta' \sqsubseteq \Delta \quad |\Delta'| < n \quad
        \Delta_{l_1} = \Delta', A \quad \Delta_{l_2} = \Delta', B}
       {\Sigma; n; R; \vec{\Delta} \vdash \kw{case}\;x\;(|\Delta'|.\,l_1)\;(|\Delta'|.\,l_2) : \langle \Theta; \Delta; O; F \rangle}
\]
A block starts in the state $\langle \emptyset; \Delta_l; \cdot;
\kw{unset} \rangle$, and blocks are acyclic as in CFG. A definition
$b_k$ of signature $(\Delta_k; A_k \rightarrow C_k)$ starts with its
environment in $\kw{rdi}$ and its argument in $\kw{rsi}$, in the state
$\langle \kw{rdi}{:}\kw{tup}(\Delta_k), \kw{rsi}{:}A_k; \emptyset; \cdot;
\kw{unset} \rangle$, and has result type $C_k$.


\section{The Graph Model}
\label{app:graph-model}

\subsection{Consistency}
\label{app:consistency}

Consistency is defined on elements, and from it on sets and lists
(Section~\ref{sec:graph-models}).  Two constants are consistent when
they are equal, two entries are consistent when consistent inputs lead
to consistent outputs, and elements that describe different kinds of
values, such as $\mathsf{left}(u)$ and $\mathsf{right}(v)$, are not
consistent. A set $D$ is consistent, written $\mathsf{con}(D)$, when
its elements are pairwise consistent, and two lists are consistent
with each other, written $V \approx V'$, when every element of one is
consistent with every element of the other.

\begin{figure}[h]
\small\figdefs
\defhead{Consistency of elements}{u \sim v}
\[
\begin{array}{@{}r@{\;}c@{\;}l@{}}
  c \sim c' & \iff & c = c'\\
  (V \mapsto w) \sim (V' \mapsto w') & \iff &
    (V \approx V' \Rightarrow w \sim w')
    \qquad
    (V \mapsto w) \sim \nu, \quad \nu \sim \nu\\
  \pairv{u}{v} \sim \pairv{u'}{v'} & \iff & u \sim u' \land v \sim v'\\
  \mathsf{tup}_i(v) \sim \mathsf{tup}_j(v') & \iff & i = j \Rightarrow v \sim v'\\
  \mathsf{left}(v) \sim \mathsf{left}(v') & \iff & v \sim v'
    \qquad
  \mathsf{right}(v) \sim \mathsf{right}(v') \iff v \sim v'
\end{array}
\]
\noindent and no other pairs of elements are consistent.

\defhead{Consistency of sets and lists}{\mathsf{con}(D) \qquad V \approx V'}
\[
\begin{array}{@{}r@{\;}c@{\;}l@{}}
  \mathsf{con}(D) & \iff & u \sim v \text{ for all } u, v \in D\\
  V \approx V' & \iff & u \sim v \text{ for all } u \in V \text{ and } v \in V'
\end{array}
\]
\caption{Consistency of elements, sets, and lists.}
\label{fig:consistency}
\end{figure}

\subsection{Depth}
\label{app:depth}

The depth of an element of the graph model is the height of its
syntax tree, not counting constants, $\nu$ and $\langle\rangle$; the depth of a finite list is the largest depth of its
elements (zero for the empty list):
\[
\begin{array}{@{}r@{\;}c@{\;}l@{}}
  \depth(V \mapsto w) & = & 1 + \max(\depth(V), \depth(w))\\
  \depth(\pairv{u}{v}) & = & 1 + \max(\depth(u), \depth(v))\\
  \depth(t(v)) & = & 1 + \depth(v)
    \quad \text{for } t \in \{\mathsf{tup}_i, \mathsf{left}, \mathsf{right}\}\\
  \depth(v) & = & 0
    \quad \text{for } v \in \{c, \nu, \langle\rangle\}\\
  \depth(V) & = & \max_{v \in V} \depth(v)
\end{array}
\]
